% ============================================================
%  REPA – Main Slides (English) with Vietnamese speaker notes
% ============================================================

% ------------------------------------------------------------------
% SLIDE 1 – TITLE
% ------------------------------------------------------------------
\begin{frame}[plain]
  \vspace{0.4cm}
  \begin{columns}[T,totalwidth=\textwidth]
    \column{0.61\textwidth}
    {\color{signal}\large Paper Presentation\par}
    \vspace{0.2cm}
    {\color{accent}\bfseries\fontsize{19}{22}\selectfont
    Representation Alignment\par
    for Generation:\par
    Training Diffusion Transformers\par
    Is Easier Than You Think\par}
    \vspace{0.28cm}
    {\large Yu et al.\hspace{0.25cm}|\hspace{0.25cm}ICLR 2025\par}
    \vspace{0.28cm}
    \begin{beamercolorbox}[wd=\linewidth,rounded=true,sep=0.28cm]{block body}
      \textbf{Talk outline.} Why diffusion transformers hit a representation bottleneck,
      how REPA turns it into an explicit alignment objective, and why this yields
      both faster convergence and state-of-the-art FID.
    \end{beamercolorbox}
    \vfill
    {\small Figures are reproduced from the original paper.}
    \column{0.37\textwidth}
    \begin{tikzpicture}[every node/.style={inner sep=0pt}]
      \node[draw=accent,fill=accentlight,rounded corners=12pt,
            minimum width=5.0cm,minimum height=6.6cm] (card) {};
      \node[anchor=north west,text width=4.4cm,align=left]
        at ($(card.north west)+(0.28cm,-0.32cm)$) {
        \textbf{Key claims}
        \begin{itemize}
          \item Better internal features make generation easier.
          \item Align early hidden states with pretrained visual semantics.
          \item SiT-XL/2 reaches strong quality in far fewer training steps.
        \end{itemize}
      };
      \node[draw=signal,fill=paperbg,rounded corners=10pt,
            minimum width=4.3cm,minimum height=1.4cm]
        at ($(card.south)+(0,1.0cm)$) {
        \begin{tabular}{c}
          \textbf{Headline result}\\
          FID \(=1.42\) on ImageNet 256\(\times\)256
        \end{tabular}
      };
    \end{tikzpicture}
  \end{columns}
  \note{
    Xin chào mọi người. Hôm nay mình sẽ trình bày về paper
    ``Representation Alignment for Generation'' hay còn gọi tắt là REPA,
    được công bố tại ICLR 2025.

    Ý tưởng chính của bài báo rất đơn giản nhưng hiệu quả:
    thay vì để diffusion transformer tự học biểu diễn ngữ nghĩa chỉ qua
    việc khử nhiễu, ta ép buộc các hidden state ở tầng sớm phải
    căn chỉnh (align) với biểu diễn từ một encoder đã được pretrain tốt
    như DINOv2.

    Kết quả đạt FID 1.42 trên ImageNet 256x256 — state-of-the-art
    tại thời điểm công bố.
  }
\end{frame}

% ------------------------------------------------------------------
% SLIDE 2 – ONE-SLIDE SUMMARY
% ------------------------------------------------------------------
\begin{frame}[t]
  \framesubtitle{Overview}
  \frametitle{One-slide summary: make representation learning explicit}
  \begin{columns}[T,totalwidth=\textwidth]
    \column{0.47\textwidth}
    \callout{\textbf{Core idea.} Instead of hoping the diffusion transformer discovers
    strong semantics through denoising alone, REPA \emph{explicitly} aligns an early
    hidden state with a clean representation from a pretrained encoder (DINOv2).}
    \vspace{0.18cm}
    \begin{itemize}
      \item Hidden states become semantically stronger early in training.
      \item Later layers can focus on high-frequency detail.
      \item The same model family gets both faster convergence and better quality.
    \end{itemize}
    \vspace{0.1cm}
    \begin{beamercolorbox}[wd=\linewidth,rounded=true,sep=0.18cm]{block body}
      \textbf{Headline.} SiT-XL/2 + REPA matches or exceeds the vanilla 7M-step
      regime in a fraction of the training budget.
    \end{beamercolorbox}
  \column{0.51\textwidth}
    \FigOne[width=\linewidth]
  \end{columns}
  \note{
    Đây là slide tổng quan. Ý tưởng cốt lõi của REPA: diffusion
    transformer khi train bình thường chỉ có tín hiệu giám sát từ bài
    toán khử nhiễu — đây là tín hiệu yếu để học biểu diễn ngữ nghĩa.

    REPA thêm một auxiliary loss để ép hidden state ở block sớm phải
    giống với feature từ DINOv2 — một encoder đã được train trên dữ liệu
    không nhãn với khả năng biểu diễn ngữ nghĩa rất tốt.

    Kết quả: SiT-XL/2 + REPA chỉ cần khoảng 400K iteration để đạt
    chất lượng mà phiên bản gốc cần 7 triệu iteration. Hình bên phải
    là Figure 1 của paper thể hiện tổng quan phương pháp.
  }
\end{frame}

% ------------------------------------------------------------------
% SLIDE 3 – PROBLEM STATEMENT
% ------------------------------------------------------------------
\begin{frame}[t]
  \framesubtitle{Motivation}
  \frametitle{Why representation quality is the bottleneck}
  \begin{columns}[T,totalwidth=\textwidth]
    \column{0.53\textwidth}
    \begin{itemize}
      \item A denoising model must recover both \textit{semantics} and
            \textit{fine detail} from noisy latents.
      \item Pure reconstruction pressure is \textbf{weak supervision} for
            learning noise-invariant, discriminative features.
      \item A large transformer wastes training budget rediscovering semantics
            that modern SSL encoders already know.
    \end{itemize}
    \vspace{0.18cm}
    \callout{\textbf{Hypothesis.} The hard part of training DiT/SiT is not
    solving the denoising objective \emph{per se} — it is learning an internal
    representation good enough to support generation.}
    \vspace{0.18cm}
    {\small Denoising gives an implicit encoder \(f_\theta(z_t)=h_t\), but
    nothing guarantees \(h_t\) will align with the semantic structure of
    strong visual encoders.}
  \column{0.43\textwidth}
    \begin{tikzpicture}[>=Latex,thick,node distance=0.55cm]
      \tikzstyle{flowbox}=[draw=accent,fill=accentlight,rounded corners=6pt,
        minimum width=4.4cm,minimum height=0.95cm,align=center]
      \node[flowbox] (noise) {Noisy latent \(z_t\)};
      \node[flowbox,below=of noise] (encoder) {DiT / SiT hidden state \(h_t\)};
      \node[flowbox,below=of encoder] (decode) {Velocity or score prediction};
      \draw[->,color=signal] (noise) -- (encoder);
      \draw[->,color=signal] (encoder) -- (decode);
      \node[draw=signal,fill=signalsoft,rounded corners=8pt,
            text width=4.2cm,align=left,below=0.8cm of decode] (issue) {
        \textbf{Issue}\\
        Semantics are learned only as a side effect of reconstruction
      };
      \draw[->,color=signal] (decode) -- (issue);
    \end{tikzpicture}
    \vspace{0.15cm}
    \begin{beamercolorbox}[wd=\linewidth,rounded=true,sep=0.18cm]{block body}
      Better semantics should reduce the optimization burden.
    \end{beamercolorbox}
  \end{columns}
  \note{
    Slide này giải thích vấn đề cốt lõi. Khi huấn luyện diffusion
    transformer, mô hình cần đồng thời học hai thứ: (1) biểu diễn
    ngữ nghĩa — hiểu ``cái gì đang có trong ảnh'', và (2) chi tiết
    tần số cao — texture, cạnh, màu sắc chính xác.

    Tín hiệu giám sát duy nhất là loss khử nhiễu — tức là bắt mô hình
    dự đoán lại noise hoặc velocity. Đây là tín hiệu yếu cho việc học
    biểu diễn ngữ nghĩa vì nó chủ yếu tập trung vào reconstruction.

    Sơ đồ bên phải cho thấy luồng: latent nhiễu đi qua transformer
    tạo ra hidden state, rồi predict velocity. Vấn đề là semantics
    chỉ được học như ``hiệu ứng phụ'' của quá trình reconstruction.

    Giả thuyết của paper: phần khó nhất không phải là giải bài toán
    denoising, mà là học được biểu diễn nội tại đủ tốt để hỗ trợ
    quá trình sinh ảnh.
  }
\end{frame}

% ------------------------------------------------------------------
% SLIDE 4 – EVIDENCE: DIFFUSION FEATURES
% ------------------------------------------------------------------
\begin{frame}[t]
  \framesubtitle{Evidence}
  \frametitle{Diffusion features are meaningful but still misaligned}
  \FigTwo[width=\linewidth]
  \vspace{0.08cm}
  \begin{columns}[T,totalwidth=\textwidth]
    \column{0.32\textwidth}
    {\small \textbf{Observation 1}\\
    SiT hidden states are not useless: linear probing peaks at later
    layers and shows non-trivial semantics.}
    \column{0.32\textwidth}
    {\small \textbf{Observation 2}\\
    Alignment to DINOv2 exists but is much weaker than alignment
    between strong SSL models.}
    \column{0.32\textwidth}
    {\small \textbf{Observation 3}\\
    Bigger models and longer training improve alignment only slowly —
    brute-force scaling is expensive.}
  \end{columns}
  \note{
    Đây là Figure 2 trong paper — bằng chứng thực nghiệm trước khi
    giới thiệu REPA. Ba quan sát chính:

    Thứ nhất, hidden state của SiT không phải vô nghĩa — khi làm
    linear probing trên ImageNet, accuracy đạt đỉnh ở các layer sâu
    hơn, cho thấy có semantic information.

    Thứ hai, tuy nhiên, mức độ alignment với DINOv2 rất yếu so với
    alignment giữa các mô hình SSL mạnh với nhau. Tức là DINOv2 và
    CLIP align tốt với nhau, nhưng SiT thì không align tốt với
    bất kỳ cái nào.

    Thứ ba, khi tăng kích thước mô hình hoặc train lâu hơn, alignment
    cải thiện rất chậm. Điều này cho thấy brute-force scaling không
    phải cách hiệu quả — ta cần một phương pháp trực tiếp hơn.
  }
\end{frame}

% ------------------------------------------------------------------
% SLIDE 5 – PRELIMINARIES: STOCHASTIC INTERPOLANTS
% ------------------------------------------------------------------
\begin{frame}[t]
  \framesubtitle{Background}
  \frametitle{Preliminaries: stochastic interpolants}
  \begin{columns}[T,totalwidth=\textwidth]
    \column{0.56\textwidth}
    \textbf{Forward interpolation}
    \begin{align*}
      x_t &= \alpha_t\, x^\ast + \sigma_t\, \epsilon,\qquad
      \epsilon \sim \mathcal{N}(0, I) \\[4pt]
      \alpha_0 &= \sigma_1 = 1,\qquad \alpha_1 = \sigma_0 = 0
    \end{align*}
    \vspace{-0.2cm}
    \begin{itemize}
      \item \(x^\ast\): clean data sample
      \item \(t \in [0,1]\) interpolates from data (\(t{=}0\)) to noise (\(t{=}1\))
      \item SiT uses the linear schedule: \(\alpha_t = 1-t,\;\sigma_t = t\)
    \end{itemize}
    \vspace{0.1cm}
    \textbf{Velocity field}
    \begin{align*}
      \dot{x}_t = v(x_t, t)
    \end{align*}
    The probability flow ODE and its reverse SDE give equivalent generation paths.
    REPA is agnostic to the exact formulation — it only needs access to hidden states.
  \column{0.40\textwidth}
    \begin{tikzpicture}[>=Latex,thick]
      \draw[line width=1pt,color=accent] (0,0) -- (4.8,0);
      \foreach \x/\label in {0/{\(t{=}0\)},2.4/{},4.8/{\(t{=}1\)}} {
        \fill[color=signal] (\x,0) circle (1.9pt);
        \node[anchor=north] at (\x,-0.1) {\label};
      }
      \node[draw=accent,fill=accentlight,rounded corners=7pt,
            minimum width=1.65cm,minimum height=1.0cm] at (0,1.1) {clean data};
      \node[draw=accent,fill=softgray,rounded corners=7pt,
            minimum width=1.75cm,minimum height=1.0cm] at (2.4,1.1) {mixed signal};
      \node[draw=signal,fill=signalsoft,rounded corners=7pt,
            minimum width=1.65cm,minimum height=1.0cm] at (4.8,1.1) {Gaussian noise};
      \draw[->,color=accent] (0.85,1.1) -- (1.55,1.1);
      \draw[->,color=accent] (3.25,1.1) -- (3.95,1.1);
      \node[anchor=west,text width=4.85cm,align=left] at (-0.05,-1.0) {
        \textbf{Key point for REPA}\\
        REPA supervises a hidden state from \emph{noisy} input,
        while the target representation comes from the \emph{clean} image.
      };
    \end{tikzpicture}
  \end{columns}
  \note{
    Trước khi đi vào REPA, cần nắm nền tảng toán học. Paper dùng
    framework stochastic interpolants — đây là cách nhìn liên tục
    về quá trình forward corruption.

    Ta nội suy giữa ảnh sạch $x^*$ và nhiễu Gaussian $\epsilon$ bằng
    hệ số $\alpha_t$ và $\sigma_t$. SiT dùng lịch trình tuyến tính
    đơn giản: $\alpha_t = 1 - t$, $\sigma_t = t$.

    Khi $t = 0$ ta có ảnh sạch, $t = 1$ là nhiễu thuần. Mô hình cần
    học velocity field $v(x_t, t)$ — tức là hướng và tốc độ di chuyển
    trên đường từ nhiễu về dữ liệu.

    Điểm quan trọng cho REPA: mô hình nhận input nhiễu nhưng target
    alignment lại đến từ ảnh sạch — đây là key design choice.
  }
\end{frame}

% ------------------------------------------------------------------
% SLIDE 6 – TRAINING OBJECTIVE
% ------------------------------------------------------------------
\begin{frame}[t]
  \framesubtitle{Background}
  \frametitle{Training objective recap}
  \begin{columns}[T,totalwidth=\textwidth]
    \column{0.51\textwidth}
    \textbf{Velocity objective}
    \begin{align*}
      v(x, t) &= \dot{\alpha}_t\,\mathbb{E}[x^\ast \mid x_t{=}x]
                 + \dot{\sigma}_t\,\mathbb{E}[\epsilon \mid x_t{=}x] \\[4pt]
      \mathcal{L}_{\text{vel}}(\theta)
      &= \mathbb{E}_{x^\ast,\epsilon,t}
      \bigl\|v_\theta(x_t, t) - \dot{\alpha}_t x^\ast
            - \dot{\sigma}_t \epsilon\bigr\|_2^2
    \end{align*}
    \vspace{0.08cm}
    \textbf{Reverse SDE}
    \begin{align*}
      dx_t = v(x_t, t)\,dt
             - \tfrac{1}{2}\,w_t\,s(x_t, t)\,dt
             + w_t\,d\bar{w}_t
    \end{align*}
  \column{0.47\textwidth}
    \textbf{Score from velocity}
    \begin{align*}
      s(x,t) = \sigma_t^{-1} \cdot
      \frac{\alpha_t v(x,t) - \dot{\alpha}_t x}
           {\dot{\alpha}_t \sigma_t - \alpha_t \dot{\sigma}_t}
    \end{align*}
    \vspace{0.16cm}
    \callout{\textbf{Interpretation.} The denoising loss teaches the network
    to invert the corruption process. REPA adds a second pressure:
    hidden states should also look like clean, semantically rich visual features.}
    \vspace{0.16cm}
    {\small Applicable to both SiT-style velocity prediction and
    DiT-style DDPM objectives.}
  \end{columns}
  \note{
    Slide này nhắc lại training objective. Mô hình học velocity field
    bằng cách minimize khoảng cách giữa velocity dự đoán và velocity
    thực tế (tổ hợp tuyến tính của $x^*$ và $\epsilon$).

    Khi sampling, ta dùng reverse SDE hoặc ODE để đi ngược từ nhiễu
    về dữ liệu. Score function $s(x,t)$ có thể được tính từ velocity
    qua công thức chuyển đổi.

    Điểm mấu chốt: loss khử nhiễu chỉ dạy mô hình ``đảo ngược quá
    trình corruption''. REPA sẽ thêm một áp lực thứ hai: hidden state
    phải trông giống feature ngữ nghĩa sạch.

    Paper áp dụng cho cả SiT (velocity) lẫn DiT (DDPM noise prediction).
  }
\end{frame}

% ------------------------------------------------------------------
% SLIDE 7 – REPA INTUITION
% ------------------------------------------------------------------
\begin{frame}[t]
  \framesubtitle{REPA}
  \frametitle{Intuition: noisy hidden states should predict clean semantics}
  \begin{columns}[T,totalwidth=\textwidth]
    \column{0.43\textwidth}
    \begin{itemize}
      \item A \textbf{pretrained encoder} sees the clean image \(x^\ast\)
            and produces target features \(y^\ast\).
      \item The \textbf{diffusion transformer} sees noisy latent \(z_t\)
            and produces hidden state \(h_t\).
      \item A small \textbf{projector MLP} maps \(h_t\) into the target
            feature space, and REPA \textbf{aligns them patch by patch}.
    \end{itemize}
    \vspace{0.14cm}
    \callout{\textbf{Elegance.} No redesign of sampling or backbone.
    Just a simple regularizer that distills external semantics into
    the denoising model during training.}
  \column{0.55\textwidth}
    \FigThree[width=\linewidth]
  \end{columns}
  \note{
    Đây là ý tưởng trực quan của REPA — slide quan trọng nhất.

    Ta có hai nhánh song song:
    — Nhánh trên: DINOv2 (frozen, pretrained) nhận ảnh sạch $x^*$
      và cho ra feature map $y^*$ với ngữ nghĩa mạnh.
    — Nhánh dưới: diffusion transformer nhận latent nhiễu $z_t$,
      đi qua các block, tạo ra hidden state $h_t$ ở một block sớm.

    Một MLP projector nhỏ (3 layer) ánh xạ $h_t$ vào không gian
    feature của DINOv2, rồi ta tính loss alignment patch-by-patch.

    Điểm đẹp: ta không thay đổi kiến trúc backbone, không thay đổi
    sampler, chỉ thêm một auxiliary loss. Projector MLP chỉ dùng khi
    train, bỏ khi inference.

    Hình bên phải là Figure 3 trong paper, minh họa pipeline này.
  }
\end{frame}

% ------------------------------------------------------------------
% SLIDE 8 – FORMAL OBJECTIVE
% ------------------------------------------------------------------
\begin{frame}[t]
  \framesubtitle{REPA}
  \frametitle{Formal objective}
  \begin{columns}[T,totalwidth=\textwidth]
    \column{0.58\textwidth}
    \textbf{Target representation from pretrained encoder:}
    \begin{align*}
      y^\ast = f(x^\ast) \in \mathbb{R}^{N \times D}
    \end{align*}
    \textbf{Hidden state from diffusion model:}
    \begin{align*}
      h_t = f_\theta(z_t) \in \mathbb{R}^{N \times D'}
    \end{align*}
    \textbf{REPA loss (patch-wise alignment):}
    \begin{align*}
      \mathcal{L}_{\text{REPA}}(\theta, \phi)
      = -\,\mathbb{E}_{x^\ast,\epsilon,t}\!\left[
        \frac{1}{N}\sum_{n=1}^{N}
        \mathrm{sim}\!\bigl(y_n^\ast,\;h_\phi(h_t)_n\bigr)
      \right]
    \end{align*}
    \textbf{Combined training objective:}
    \begin{align*}
      \boxed{\;\mathcal{L} = \mathcal{L}_{\text{velocity}}
        + \lambda\,\mathcal{L}_{\text{REPA}}\;}
    \end{align*}
    \vspace{0.08cm}
    \begin{itemize}
      \item \(f\): frozen pretrained encoder (typically DINOv2)
      \item \(h_\phi\): small trainable projection head (3-layer MLP)
      \item \(\mathrm{sim}\): NT-Xent or cosine similarity
      \item \(\lambda\): tradeoff weight (default \(\approx 0.5\))
    \end{itemize}
  \column{0.38\textwidth}
    \begin{beamercolorbox}[wd=\linewidth,rounded=true,sep=0.20cm]{block body}
      \textbf{Key design choice}\\[4pt]
      The target comes from the \emph{clean image}, not from a noisy
      augmentation. The model learns to recover semantic structure
      that is invariant to the diffusion corruption.
    \end{beamercolorbox}
    \vspace{0.18cm}
    \begin{beamercolorbox}[wd=\linewidth,rounded=true,sep=0.20cm]{block body}
      \textbf{Operationally}\\[4pt]
      REPA is just an auxiliary loss. It slots into existing DiT/SiT
      training with minimal code change.
    \end{beamercolorbox}
    \vspace{0.18cm}
    \begin{beamercolorbox}[wd=\linewidth,rounded=true,sep=0.20cm]{block body}
      \textbf{Notation}\\[4pt]
      \(N\) = number of spatial patches\\
      \(D, D'\) = feature dimensions\\
      \(n\) indexes individual patches
    \end{beamercolorbox}
  \end{columns}
  \note{
    Slide công thức chính thức — hãy giải thích từng phần.

    Đầu tiên, encoder pretrained $f$ (DINOv2) nhận ảnh sạch $x^*$
    và cho ra ma trận feature $y^* \in \mathbb{R}^{N \times D}$,
    trong đó $N$ là số patch (ví dụ 256 patch cho ảnh 256x256 qua
    VAE thành 32x32, rồi patchify), $D$ là chiều feature.

    Hidden state $h_t$ từ diffusion model có chiều $D'$ khác, nên
    cần projection head $h_\phi$ (MLP 3 layer) để ánh xạ về cùng
    không gian.

    REPA loss tính similarity giữa từng cặp patch tương ứng, rồi
    lấy trung bình trên $N$ patch. Similarity có thể là cosine
    hoặc NT-Xent (contrastive).

    Loss tổng hợp: cộng velocity loss với $\lambda$ nhân REPA loss.
    $\lambda$ mặc định khoảng 0.5 — paper cho thấy kết quả khá
    robust với giá trị này.

    Lưu ý quan trọng: target đến từ ảnh SẠCH, không phải ảnh nhiễu.
    Đây là design choice then chốt — model phải học recover semantics
    bất biến với corruption.
  }
\end{frame}

% ------------------------------------------------------------------
% SLIDE 9 – ARCHITECTURE: WHERE REPA ATTACHES
% ------------------------------------------------------------------
\begin{frame}[t]
  \framesubtitle{REPA}
  \frametitle{Architecture: where REPA is attached}
  \begin{columns}[T,totalwidth=\textwidth]
    \column{0.46\textwidth}
    \FigNine[width=\linewidth]
  \column{0.50\textwidth}
    \begin{itemize}
      \item Backbone is identical to DiT/SiT: patchify, run transformer blocks,
            predict velocity.
      \item REPA taps an \textbf{intermediate hidden state},
            typically from an \textbf{early block} (depth 6--8).
      \item A lightweight \textbf{3-layer MLP projector} (with SiLU) bridges
            the hidden state to the target feature space.
      \item The projector is \textbf{training-time only} —
            discarded at inference.
      \item AdaLN-zero conditioning and the denoising head
            remain \textbf{untouched}.
    \end{itemize}
    \vspace{0.15cm}
    \callout{\textbf{Takeaway.} REPA is not a new architecture. It is a
    training-time regularization layer on top of the existing model.}
  \end{columns}
  \note{
    Slide kiến trúc. Hình bên trái trích từ paper (appendix Figure 9)
    cho thấy cách REPA được gắn vào mô hình.

    Backbone hoàn toàn giữ nguyên: ảnh latent được patchify thành
    token sequence, đi qua các transformer block với AdaLN-zero
    conditioning (inject timestep $t$ và class label).

    REPA ``tap'' vào hidden state ở một block trung gian sớm —
    thường là block thứ 6 hoặc 8 (trong tổng số 28 block của XL).

    Projection head là MLP 3 layer với activation SiLU, rất nhẹ.
    Nó chỉ dùng khi training để ánh xạ hidden state vào không gian
    DINOv2. Khi inference, ta bỏ hoàn toàn projector — model chạy
    giống hệt SiT/DiT gốc.

    Đây là điểm mạnh: REPA không thay đổi kiến trúc hay sampler,
    chỉ thêm auxiliary loss, nên rất dễ tích hợp.
  }
\end{frame}

% ------------------------------------------------------------------
% SLIDE 10 – WHY EARLY LAYERS
% ------------------------------------------------------------------
\begin{frame}[t]
  \framesubtitle{REPA}
  \frametitle{Why early-layer alignment is enough}
  \begin{columns}[T,totalwidth=\textwidth]
    \column{0.49\textwidth}
    \begin{tabular}{L{2.2cm}C{1.1cm}C{1.2cm}L{2.4cm}}
      \toprule
      Setting & Depth & FID \metricdown & Interpretation \\
      \midrule
      Vanilla SiT-L/2 & -- & 18.8 & no alignment \\
      REPA + DINOv2-L & 6 & 10.3 & strong early gains \\
      REPA + DINOv2-L & 8 & \textbf{10.0} & best tradeoff \\
      REPA + DINOv2-L & 12 & 11.2 & alignment hurts \\
      REPA + DINOv2-L & 16 & 12.1 & too constrained \\
      \bottomrule
    \end{tabular}
    \vspace{0.18cm}
    \papersource{Table 3 in the paper}
  \column{0.47\textwidth}
    \begin{enumerate}
      \item Early blocks learn a \textbf{semantic scaffold}.
      \item Middle and late blocks no longer need to rediscover
            object-level structure.
      \item Remaining capacity focuses on \textbf{high-frequency,
            generative detail}.
    \end{enumerate}
    \vspace{0.18cm}
    \callout{\textbf{Interpretation.} REPA should shape the
    representation space \emph{early}, then step out of the way
    so the model can finish the denoising task with later blocks.}
  \end{columns}
  \note{
    Tại sao align ở layer sớm lại hiệu quả? Bảng trái cho thấy
    kết quả thí nghiệm thay đổi vị trí alignment depth.

    Depth 6: FID giảm mạnh từ 18.8 xuống 10.3 — cải thiện lớn.
    Depth 8: tốt nhất, đạt FID 10.0.
    Depth 12 và 16: bắt đầu xấu đi — alignment ở quá sâu
    ``chèn ép'' khả năng của các layer sau.

    Giải thích: các block sớm nên lo việc xây dựng ``khung ngữ nghĩa''
    — hiểu scene có gì, vật thể ở đâu. Các block giữa và sau tập
    trung vào chi tiết tần số cao — texture, edge, color chính xác.

    Nếu ta ép alignment ở quá nhiều layer hoặc layer sâu, model
    không còn đủ ``dư địa'' cho phần generative detail.

    Đây là discovery thực nghiệm — paper thừa nhận chưa có lý thuyết
    chặt chẽ giải thích, nhưng kết quả rất nhất quán.
  }
\end{frame}

% ------------------------------------------------------------------
% SLIDE 11 – COMPONENT ABLATION (Table 2)
% ------------------------------------------------------------------
\begin{frame}[t]
  \framesubtitle{Experiments}
  \frametitle{Ablation: which components matter most}
  \begin{columns}[T,totalwidth=\textwidth]
    \column{0.54\textwidth}
    \begin{tabular}{L{3.2cm}C{1.1cm}C{1.2cm}C{1.1cm}}
      \toprule
      Setting (SiT-L/2, 400K) & FID \metricdown & IS \metricup & Acc \metricup \\
      \midrule
      Vanilla & 18.8 & 72.0 & -- \\
      + MAE-L & 12.5 & 90.7 & 57.3 \\
      + CLIP-L & 11.0 & 100.4 & 67.2 \\
      + SigLIP-L & 10.2 & 107.0 & 68.8 \\
      + DINOv2-L & 10.0 & 106.6 & 68.1 \\
      + DINOv2-B & 9.7 & 107.5 & 65.7 \\
      + cosine sim & 9.9 & 111.9 & 68.2 \\
      \bottomrule
    \end{tabular}
    \vspace{0.14cm}
    \papersource{Table 2 in the paper}
  \column{0.42\textwidth}
    \begin{itemize}
      \item Stronger target encoders improve both semantics and generation.
      \item Very large encoders are \emph{not} mandatory; DINOv2-B already works well.
      \item Simple cosine similarity is competitive with contrastive (NT-Xent) objectives.
      \item \textbf{Every} design choice beats the vanilla baseline by a wide margin.
    \end{itemize}
    \vspace{0.12cm}
    \callout{REPA is \textbf{robust}: the improvement is not tied to one specific
    target encoder or similarity function.}
  \end{columns}
  \note{
    Slide ablation — rất quan trọng để hiểu yếu tố nào đóng góp.

    Vanilla SiT-L/2 sau 400K iteration: FID 18.8. Khi thêm REPA:
    — MAE-L: giảm xuống 12.5. MAE học reconstruction nên features
      không mạnh về discriminative.
    — CLIP-L: 11.0 — tốt hơn nhờ features discriminative mạnh
      (contrastive image-text).
    — SigLIP-L: 10.2 — tương tự CLIP nhưng dùng sigmoid loss.
    — DINOv2-L: 10.0 — self-supervised, features rất tốt cho
      dense prediction.
    — DINOv2-B (nhỏ hơn): 9.7 — bất ngờ là nhỏ hơn lại tốt hơn
      một chút, có thể do matching dimension dễ hơn.
    — Dùng cosine similarity thay NT-Xent: 9.9 — gần như tương đương.

    Kết luận: REPA robust — không phụ thuộc vào một encoder hay
    similarity function cụ thể. Bất kỳ encoder pretrained tốt nào
    cũng cải thiện đáng kể.
  }
\end{frame}

% ------------------------------------------------------------------
% SLIDE 12 – SCALABILITY
% ------------------------------------------------------------------
\begin{frame}[t]
  \framesubtitle{Experiments}
  \frametitle{Scalability: stronger targets and larger DiTs help more}
  \FigFive[width=\linewidth]
  \vspace{0.05cm}
  \begin{columns}[T,totalwidth=\textwidth]
    \column{0.32\textwidth}
    {\small \textbf{Target quality}\\
    Better pretrained encoders push both linear-probe accuracy
    and FID in the right direction.}
    \column{0.32\textwidth}
    {\small \textbf{Model size}\\
    Larger SiT models benefit more from REPA; the relative
    speedup grows with scale.}
    \column{0.32\textwidth}
    {\small \textbf{Joint frontier}\\
    REPA improves the discrimination--generation frontier rather
    than trading one metric for another.}
  \end{columns}
  \note{
    Figure 5 cho thấy tính scalability. Ba điểm chính:

    Thứ nhất — target quality: encoder tốt hơn (DINOv2 > CLIP > MAE)
    thì cả linear probe accuracy lẫn FID đều cải thiện. Tức là chất
    lượng representation target trực tiếp ảnh hưởng kết quả.

    Thứ hai — model size: SiT-B, SiT-L, SiT-XL — model lớn hơn
    hưởng lợi nhiều hơn từ REPA. Relative speedup tăng theo scale.

    Thứ ba — REPA cải thiện cả hai chiều: discrimination (accuracy)
    lẫn generation (FID) cùng lúc, thay vì đánh đổi một cái cho
    cái kia. Đây là frontier improvement thực sự.
  }
\end{frame}

% ------------------------------------------------------------------
% SLIDE 13 – CONVERGENCE
% ------------------------------------------------------------------
\begin{frame}[t]
  \framesubtitle{Experiments}
  \frametitle{Convergence and efficiency}
  \begin{columns}[T,totalwidth=\textwidth]
    \column{0.48\textwidth}
    \FigEleven[width=\linewidth]
  \column{0.48\textwidth}
    \begin{tabular}{L{2.3cm}C{1.2cm}C{1.1cm}}
      \toprule
      Model & Iterations & FID \metricdown \\
      \midrule
      SiT-XL/2 & 7M & 8.3 \\
      \midrule
      REPA & 400K & 7.9 \\
      REPA & 1M & 6.4 \\
      REPA & 4M & 5.9 \\
      \midrule
      DiT-XL/2 & 7M & 9.6 \\
      REPA & 850K & 9.6 \\
      \bottomrule
    \end{tabular}
    \vspace{0.14cm}
    \callout{\textbf{Key message.} REPA is not only a better endpoint:
    it changes the optimization trajectory so the model enters a good
    semantic regime much earlier.}
    \vspace{0.12cm}
    {\scriptsize The paper highlights a \(\mathbf{>17.5\times}\) speedup for
    matching strong SiT performance without CFG.}
  \end{columns}
  \note{
    Slide tốc độ hội tụ — một trong những kết quả ấn tượng nhất.

    Bảng bên phải so sánh trực tiếp:
    — SiT-XL/2 vanilla: 7 triệu iteration mới đạt FID 8.3.
    — REPA: chỉ 400K iteration đã đạt FID 7.9 — tốt hơn!
      Tức là nhanh hơn hơn 17.5 lần.
    — Train thêm đến 4M iteration: FID giảm xuống 5.9.

    Với DiT-XL/2: vanilla cần 7M iteration cho FID 9.6,
    REPA chỉ cần 850K — nhanh hơn 8 lần.

    Hình bên trái (Figure 11) show training curve — đường REPA
    giảm nhanh hơn rất nhiều ngay từ đầu training.

    Insight quan trọng: REPA không chỉ cải thiện kết quả cuối cùng,
    mà thay đổi toàn bộ trajectory tối ưu hóa. Model ``vào vùng
    ngữ nghĩa tốt'' sớm hơn nhiều.
  }
\end{frame}

% ------------------------------------------------------------------
% SLIDE 14 – SYSTEM-LEVEL COMPARISON
% ------------------------------------------------------------------
\begin{frame}[t]
  \framesubtitle{Experiments}
  \frametitle{System-level comparison on ImageNet 256\(\times\)256}
  \begin{columns}[T,totalwidth=\textwidth]
    \column{0.58\textwidth}
    \begin{tabular}{L{3.4cm}C{1.1cm}C{1.2cm}C{1.0cm}}
      \toprule
      Method & Epochs & FID \metricdown & Rec. \metricup \\
      \midrule
      U-ViT-H/2 & 240 & 2.29 & 0.57 \\
      DiffiT* & -- & 1.73 & 0.62 \\
      MDTv2-XL/2* & 1080 & 1.58 & 0.65 \\
      \midrule
      DiT-XL/2 & 1400 & 2.27 & 0.57 \\
      SiT-XL/2 & 1400 & 2.06 & 0.59 \\
      \midrule
      REPA (800 ep) & 800 & 1.80 & 0.61 \\
      REPA + guid.\ interval* & 800 & \textbf{1.42} & \textbf{0.64} \\
      \bottomrule
    \end{tabular}
    \vspace{0.08cm}
    {\scriptsize * Uses guidance interval or advanced CFG scheduling.}
  \column{0.38\textwidth}
    \begin{itemize}
      \item REPA beats vanilla SiT with fewer epochs.
      \item With CFG and guidance interval: \textbf{FID = 1.42} — best
            in class at time of publication.
      \item Competitive against strong transformer and hybrid baselines.
    \end{itemize}
    \vspace{0.16cm}
    \callout{\textbf{Bottom line.} REPA + guidance interval is
    the best performing DiT-family model on ImageNet 256\(\times\)256
    at publication time.}
  \end{columns}
  \note{
    So sánh system-level trên ImageNet 256x256 — benchmark chuẩn.

    Các baseline mạnh: DiT-XL/2 cần 1400 epoch để đạt FID 2.27,
    SiT-XL/2 cần 1400 epoch cho FID 2.06.

    REPA chỉ cần 800 epoch (gần một nửa) đã đạt FID 1.80 — tốt
    hơn đáng kể. Khi kết hợp với guidance interval (một kỹ thuật
    CFG scheduling), đạt FID 1.42 — state-of-the-art.

    So với các phương pháp khác: MDTv2-XL/2 đạt 1.58 nhưng cần
    1080 epoch và dùng kiến trúc khác. DiffiT đạt 1.73. REPA
    vượt trội tất cả.

    Recall ratio cũng tốt: 0.64 — cho thấy model không chỉ
    generate ảnh giống training data mà còn đa dạng.
  }
\end{frame}

% ------------------------------------------------------------------
% SLIDE 15 – QUALITATIVE
% ------------------------------------------------------------------
\begin{frame}[t]
  \framesubtitle{Experiments}
  \frametitle{Qualitative improvement appears early in training}
  \FigFour[width=\linewidth]
  \vspace{0.06cm}
  \begin{columns}[T,totalwidth=\textwidth]
    \column{0.32\textwidth}
    {\small \textbf{Controlled comparison}\\
    Same noise seed, same sampler, same number of function evaluations, no CFG.}
    \column{0.32\textwidth}
    {\small \textbf{Visual pattern}\\
    REPA sharpens object identity and compositional structure earlier.}
    \column{0.32\textwidth}
    {\small \textbf{Interpretation}\\
    Better internal semantics translate into visibly better
    generations during the early training phase.}
  \end{columns}
  \note{
    Kết quả trực quan từ Figure 4. So sánh controlled: cùng noise
    seed, cùng sampler, cùng số bước sampling, không dùng CFG.

    Hàng trên: SiT vanilla ở các mốc training khác nhau.
    Hàng dưới: SiT + REPA ở cùng các mốc đó.

    Rõ ràng REPA cho ảnh sắc nét hơn, nhận diện vật thể tốt hơn,
    cấu trúc ảnh hợp lý hơn — đặc biệt ở giai đoạn đầu training.

    Ví dụ: ở 200K iteration, vanilla vẫn còn rất mờ và thiếu cấu
    trúc, trong khi REPA đã cho thấy vật thể rõ ràng.

    Điều này nhất quán với giả thuyết: khi semantics tốt hơn sớm,
    model có thể tập trung vào chi tiết sớm hơn.
  }
\end{frame}

% ------------------------------------------------------------------
% SLIDE 16 – WHY REPA HELPS
% ------------------------------------------------------------------
\begin{frame}[t]
  \framesubtitle{Interpretation}
  \frametitle{Why REPA helps both speed and final FID}
  \begin{center}
  \begin{tikzpicture}[>=Latex,thick,node distance=1.15cm]
    \tikzstyle{ibox}=[draw=accent,fill=accentlight,rounded corners=8pt,
      minimum width=4.0cm,minimum height=1.1cm,align=center]
    \tikzstyle{gbox}=[draw=signal,fill=signalsoft,rounded corners=8pt,
      minimum width=4.2cm,minimum height=1.15cm,align=center]
    \node[ibox] (a) {1.\ Early hidden states\\absorb pretrained semantics};
    \node[gbox,right=1.1cm of a] (b) {2.\ Optimization becomes easier\\
      because semantics are \emph{explicit}};
    \node[ibox,right=1.1cm of b] (c) {3.\ Later layers focus on\\
      texture, geometry, detail};
    \draw[->,color=accent,line width=1.1pt] (a) -- (b);
    \draw[->,color=accent,line width=1.1pt] (b) -- (c);
  \end{tikzpicture}
  \end{center}
  \vspace{0.22cm}
  \begin{columns}[T,totalwidth=\textwidth]
    \column{0.48\textwidth}
    \callout{\textbf{Optimization story.} REPA supplies a clean semantic
    target that the denoising loss alone would need many more updates
    to discover.}
    \column{0.48\textwidth}
    \callout{\textbf{Quality story.} Once semantics are stabilized,
    the model allocates more effective capacity to the truly generative
    part of the task.}
  \end{columns}
  \vspace{0.16cm}
  {\small This interpretation is empirical, not theoretically proved,
  but consistent with all ablations in the paper.}
  \note{
    Slide giải thích ``tại sao REPA hoạt động'' — tổng hợp insight.

    Ba bước logic:

    Bước 1: Nhờ REPA, các block sớm nhanh chóng hấp thụ ngữ nghĩa
    từ DINOv2. Hidden state ở block 6-8 trở nên ``có ý nghĩa'' sớm
    trong quá trình training.

    Bước 2: Khi ngữ nghĩa đã explicit (không cần model tự khám phá),
    quá trình tối ưu hóa trở nên dễ dàng hơn. Gradient signal rõ
    ràng hơn, loss landscape smooth hơn.

    Bước 3: Các layer sau không cần lo ``understanding what's in the
    image'' nữa — chúng tập trung toàn bộ capacity vào chi tiết
    generative: texture, geometry, lighting, fine detail.

    Hai góc nhìn bổ sung: optimization story (nhanh hơn) và quality
    story (tốt hơn). Cả hai đều được paper chứng minh thực nghiệm.

    Lưu ý: đây là giải thích thực nghiệm, chưa có chứng minh lý
    thuyết chặt chẽ — paper thừa nhận đây là limitation.
  }
\end{frame}

% ------------------------------------------------------------------
% SLIDE 17 – LIMITATIONS & FUTURE
% ------------------------------------------------------------------
\begin{frame}[t]
  \framesubtitle{Discussion}
  \frametitle{Limitations and future directions}
  \begin{columns}[T,totalwidth=\textwidth]
    \column{0.52\textwidth}
    \textbf{Current limitations}
    \begin{itemize}
      \item \textbf{Alignment depth} chosen empirically; no theoretical
            guidance on why early layers are best.
      \item Mostly validated on \textbf{latent image diffusion};
            video, pixel-space, and audio generation remain open.
      \item The work is largely \textbf{empirical}; no theoretical bound
            on how much alignment helps.
      \item \textbf{Time-varying REPA weights} and richer schedules
            are left for future work.
      \item Only tested with \textbf{vision} SSL encoders; cross-modal
            targets unexplored.
    \end{itemize}
  \column{0.44\textwidth}
    \begin{beamercolorbox}[wd=\linewidth,rounded=true,sep=0.22cm]{block body}
      \textbf{Promising directions}
      \begin{itemize}
        \item Larger-scale text-to-image training
        \item Principled alignment schedules
              (vary \(\lambda\) or depth over time)
        \item Theory linking denoising and
              representation learning
        \item Broader encoder choices beyond
              current SSL vision models
        \item Video and multi-modal diffusion
      \end{itemize}
    \end{beamercolorbox}
    \vspace{0.18cm}
    \callout{REPA opens a clear line of research: what is the
    \emph{optimal} way to inject pretrained knowledge into
    generative training?}
  \end{columns}
  \note{
    Limitations — cần trung thực về những gì paper chưa giải quyết.

    Thứ nhất, alignment depth (block 6-8) được chọn qua grid search,
    không có lý thuyết giải thích tại sao. Với architecture khác
    hoặc dataset khác, con số này có thể thay đổi.

    Thứ hai, hầu hết thí nghiệm trên latent image diffusion (ImageNet).
    Video diffusion, pixel-space diffusion, audio — chưa được kiểm chứng.

    Thứ ba, hoàn toàn thực nghiệm — không có bound lý thuyết.

    Thứ tư, $\lambda$ cố định suốt training. Có thể schedule $\lambda$
    theo thời gian sẽ tốt hơn — ví dụ giảm dần khi model đã học đủ
    semantics — nhưng paper chưa explore.

    Hướng tương lai hứa hẹn: text-to-image scale lớn, lý thuyết
    liên kết denoising và representation learning, và mở rộng sang
    các modality khác.
  }
\end{frame}

% ------------------------------------------------------------------
% SLIDE 18 – TAKEAWAYS
% ------------------------------------------------------------------
\begin{frame}[t]
  \framesubtitle{Conclusion}
  \frametitle{Three key takeaways}
  \begin{columns}[T,totalwidth=\textwidth]
    \column{0.58\textwidth}
    \begin{enumerate}
      \item Diffusion transformers already learn useful features, but
            they are \textbf{not aligned enough} with strong semantic
            representations.
      \item REPA converts that weakness into a \textbf{simple auxiliary
            objective}: align noisy hidden states to clean pretrained
            features.
      \item The result is one of the clearest cases where \textbf{better
            representation learning directly improves generative
            training efficiency}.
    \end{enumerate}
    \vspace{0.18cm}
    \callout{\textbf{Bottom line.} REPA is simple, modular, and
    high-leverage: it changes neither the sampler nor the backbone,
    yet materially transforms the training dynamics.}
  \column{0.38\textwidth}
    \begin{beamercolorbox}[wd=\linewidth,rounded=true,sep=0.24cm]{block body}
      \textbf{Numbers to remember}
      \begin{itemize}
        \item FID 1.42 on ImageNet 256
        \item \(>17.5\times\) training speedup
        \item Just an auxiliary loss, minimal code change
      \end{itemize}
    \end{beamercolorbox}
    \vspace{0.24cm}
    \begin{beamercolorbox}[wd=\linewidth,rounded=true,sep=0.24cm]{block body}
      \textbf{Backup slides available}
      \begin{itemize}
        \item Derivations
        \item Metrics (CKNNA, CKA)
        \item Extra ablations
        \item 512\(\times\)512 and T2I extensions
      \end{itemize}
    \end{beamercolorbox}
    \vspace{0.24cm}
    {\small S.\ Yu et al., ``Representation Alignment for Generation,''
    ICLR 2025.}
  \end{columns}
  \note{
    Slide tổng kết — ba điểm cần nhớ:

    Một: Diffusion transformer có học features, nhưng chúng chưa đủ
    aligned với biểu diễn ngữ nghĩa mạnh. Gap này là bottleneck
    chính cho tốc độ training và chất lượng.

    Hai: REPA biến weakness thành strength bằng một auxiliary loss
    đơn giản — align hidden state với DINOv2 features. Không cần
    thay đổi architecture hay sampler.

    Ba: Đây là một trong những ví dụ rõ ràng nhất cho thấy
    ``representation learning tốt hơn = generative training hiệu
    quả hơn''. FID 1.42 và speedup 17.5 lần là con số ấn tượng.

    Nếu có câu hỏi, mình có backup slides trong appendix bao gồm
    derivations, metrics, ablations chi tiết, và extensions sang
    512x512 và text-to-image.

    Cảm ơn mọi người đã lắng nghe!
  }
\end{frame}
